\documentclass{beamer}

\usetheme{Antibes}

% Loading preambulo
% PREAMBULO ====================================================================

% Packages ---------------------------------------------------------------------

% Package to Portuguese language
\usepackage[brazilian]{babel}
% Packages to Figures
\usepackage{graphicx}
\usepackage{tikz}
% Packages to math symbols and expressions
\usepackage{amsfonts, amssymb, amsmath}
% Packages to insert code
\usepackage{listings} 
\usepackage{verbatim}
% Package to justify text
\usepackage[document]{ragged2e}
% Package to manage the bibliography
\usepackage[backend=biber, style=numeric, sorting=none]{biblatex}
% Package to facilitate quotations
\usepackage{csquotes}
% Package to use multicols
\usepackage{multicol}
% Packages for url
%	The package libertinus already loads hyperref
%\usepackage[colorlinks=true]{hyperref}
% Font
%   Fira
%\usepackage[sfdefault, lf]{FiraSans}
%   Libertinus
\usepackage{libertinus}
\renewcommand*\familydefault{\sfdefault} %% Only if the base font of the document is to be sans serif
% Colors
\usepackage{xcolor}
\usepackage{colortbl}
% Table
\usepackage{tabularray}
% Fill line
\usepackage{xhfill}

% Configurations ---------------------------------------------------------------

\AtBeginSection{
    \begin{frame}{\secname}
        \tableofcontents[currentsection,hideallsubsections]
    \end{frame}
}

\AtBeginSubsection{
    \begin{frame}{\subsecname}
        \tableofcontents[currentsection,subsectionstyle=show/shaded/hide, subsubsectionstyle=hide]
    \end{frame}
}

\AtBeginSubsubsection{
    \begin{frame}{\subsubsecname}
        \tableofcontents[currentsection,subsectionstyle=show/shaded/hide,subsubsectionstyle=show/shaded/hide/hide]
    \end{frame}
}

% Numbering slides
\setbeamertemplate{footline}[frame number]{}

% Getting rid of bottom navigation bars
\setbeamertemplate{navigation symbols}{}

% Margins
\setbeamersize{text margin left=20pt, text margin right=20pt}

% Colors
\definecolor{burgundy}{RGB}{133,24,42}
\definecolor{bloodred}{RGB}{147,5,5}

\setbeamercolor{structure}{fg=burgundy}

\setbeamercolor*{palette primary}{use=structure,fg=white,bg=structure.fg}
%\setbeamercolor*{palette secondary}{use=structure,fg=white,bg=structure.fg!75!black}
%\setbeamercolor*{palette tertiary}{use=structure,fg=white,bg=structure.fg!50!black}
\setbeamercolor*{palette quaternary}{fg=white,bg=bloodred}

%	hyperref
\hypersetup{
	colorlinks=true,
	urlcolor=blue,
	breaklinks=true}

% New commands -----------------------------------------------------------------

\NewTblrTableCommand \aula{\SetCell{bg=green!60,fg=white}}
\NewTblrTableCommand \prova{\SetCell{bg=red!80,fg=white}}
\NewTblrTableCommand \feriado{\SetCell{bg=blue!50,fg=white}}

\newcommand{\esta}[1]{\textbf{\underline{#1}}}

% Title
\title[GP 05]{Gerência de Projetos}
% Subtitle
\subtitle{Revisão AV1}
% Author of the presentation
\author[E.J.R. Silva]{Evandro J.R. Silva}
% date of the presentation
\date{}

\begin{document}
	
	\begin{frame}
		\titlepage
	\end{frame}
	
	\begin{frame}{Sumário}
		\tableofcontents
	\end{frame}

%===============================================================================
% SECTION 1 ====================================================================
%===============================================================================
\section{Princípios do Gerenciamento de Projetos}

\begin{frame}{Princípios do Gerenciamento de Projetos}
	\begin{enumerate}
		\justifying
		\item No contexto do PMBOK 7, o conceito de 'Stewardship' (Administração) refere-se a:
		\begin{enumerate}[a]
			\item A prática de microgerenciar cada tarefa técnica da equipe.
			\item \alert<2>{A responsabilidade de ser um guardião confiável dos recursos e agir de forma ética e zelosa.}
			\item O ato de delegar toda a responsabilidade financeira ao departamento de contabilidade.
			\item A aplicação estrita de metodologias ágeis em projetos de construção civil.
			\item O foco exclusivo no lucro imediato da organização, ignorando impactos sociais.
		\end{enumerate}
	\end{enumerate}
\end{frame}

\begin{frame}{Princípios do Gerenciamento de Projetos}
	\begin{enumerate}
		\setcounter{enumi}{1}
		\justifying
		\item A administração envolve quatro áreas principais: \rule{60pt}{1pt}, \rule{60pt}{1pt}, \rule{60pt}{1pt} e conformidade.
		\begin{enumerate}[a]
			\item Integridade, Sustentabilidade, Transparência
			\item \alert<2>{Integridade / Cuidado / Confiabilidade}
			\item Ética, Governança, Lucro
			\item Confiança, Liderança, Controle
			\item Responsabilidade, Autoridade, Prestação de Contas
		\end{enumerate}
	\end{enumerate}
\end{frame}

\begin{frame}{Princípios do Gerenciamento de Projetos}
	\begin{enumerate}
		\setcounter{enumi}{2}
		\justifying
		\item Quais comportamentos demonstram uma administração (stewardship) eficaz?
		\begin{enumerate}[I]
			\item Operar com transparência e honestidade.
			\item Proteger os recursos financeiros contra desperdícios.
			\item Ignorar normas de segurança para acelerar a entrega.
			\item Promover a diversidade e a inclusão na equipe.
		\end{enumerate}
		\begin{enumerate}[a]
			\item Se apenas I e II estiverem corretas.
			\item Se apenas III estiver correta.
			\item \alert<2>{Se apenas I, II e IV estiverem corretas.}
			\item Se apenas IV estiver correta.
			\item Se todas estiverem corretas.
		\end{enumerate}
	\end{enumerate}
\end{frame}

\begin{frame}{Princípios do Gerenciamento de Projetos}
	\begin{enumerate}
		\setcounter{enumi}{3}
		\justifying
		\item Relacione os pilares da administração às suas definições:
		\begin{columns}
			\begin{column}{0.28\textwidth}
				\begin{enumerate}[A]
					\item Integridade
					\item Cuidado
					\item Confiabilidade
					\item Conformidade
				\end{enumerate}
			\end{column}
			\begin{column}{0.72\textwidth}
				\begin{itemize}
					\item Agir de forma honesta e ética em todas as interações.
					\item Zelar pelo uso eficiente e sustentável dos recursos.
					\item Representar a organização de forma honrosa e cumprir compromissos.
					\item Seguir leis, regulamentos e diretrizes organizacionais.
				\end{itemize}
			\end{column}
		\end{columns}
		\begin{enumerate}[a]
			\item A, B, C, D.
			\item \alert<2>{B, C, D, A}
			\item C, D, A, B.
			\item D, A, B, C.
			\item A, C, B, D.
		\end{enumerate}
	\end{enumerate}
\end{frame}

\begin{frame}{Princípios do Gerenciamento de Projetos}
	\begin{enumerate}
		\setcounter{enumi}{4}
		\justifying
		\item A responsabilidade de 'Stewardship' no PMBOK 7 é atribuída a:
		\begin{enumerate}[a]
			\item Apenas ao patrocinador do projeto.
			\item Exclusivamente ao gerente de projeto.
			\item \alert<2>{A todos os membros da equipe e envolvidos no projeto.}
			\item Somente ao departamento jurídico.
			\item Aos stakeholders externos, como o governo.
		\end{enumerate}
	\end{enumerate}
\end{frame}

\begin{frame}{Princípios do Gerenciamento de Projetos}
	\begin{enumerate}
		\setcounter{enumi}{5}
		\justifying
		\item Ao escolher um fornecedor, você percebe que a opção mais barata utiliza mão de obra em condições precárias. Para honrar o princípio da Administração, você deve:
		\begin{enumerate}[a]
			\item Escolher o fornecedor mais barato para salvar o orçamento do projeto.
			\item Ignorar o fato, pois o que importa é o resultado técnico.
			\item \alert<2>{Recusar o fornecedor com base em critérios de ética e responsabilidade social.}
			\item Pedir um desconto ainda maior para compensar o risco moral.
			\item Deixar que o departamento de compras decida sozinho sem interferir.
		\end{enumerate}
	\end{enumerate}
\end{frame}

\begin{frame}{Princípios do Gerenciamento de Projetos}
	\begin{enumerate}
		\setcounter{enumi}{6}
		\justifying
		\item Quais ativos devem ser protegidos por um administrador zeloso?
		\begin{enumerate}[I]
			\item Ativos Financeiros (Orçamento).
			\item Ativos Físicos (Equipamentos).
			\item Capital Intelectual (Propriedade Intelectual).
			\item Reputação da Organização.
		\end{enumerate}
		\begin{enumerate}[a]
			\item Se apenas I e II estiverem corretas.
			\item Se apenas III e IV estiverem corretas.
			\item \alert<2>{Se todas estiverem corretas.}
			\item Se apenas I, II e IV estiverem corretas.
			\item Se apenas IV estiver correta.
		\end{enumerate}
	\end{enumerate}
\end{frame}

\begin{frame}{Princípios do Gerenciamento de Projetos}
	\begin{enumerate}
		\setcounter{enumi}{7}
		\justifying
		\item A criação de um ambiente colaborativo exige que a equipe desenvolva a \rule{90pt}{1pt}, que é a crença de que um membro não será punido ou humilhado por admitir erros, fazer perguntas ou propor ideias arriscadas.
		\begin{enumerate}[a]
			\item Cultura de alta performance
			\item \alert<2>{Segurança Psicológica}
			\item Hierarquia forte
			\item Responsabilidade individual
			\item Governança rígida
		\end{enumerate}
	\end{enumerate}
\end{frame}

\begin{frame}{Princípios do Gerenciamento de Projetos}
	\begin{enumerate}
		\setcounter{enumi}{8}
		\justifying
		\item Quais fatores influenciam a criação de um ambiente de equipe colaborativo?
		\begin{enumerate}[I]
			\item Estrutura da equipe (localizada ou virtual).
			\item Autoridade e papéis claramente definidos.
			\item Processos de comunicação eficazes.
			\item O uso de um único estilo de liderança para todos os membros.
		\end{enumerate}
		\begin{enumerate}[a]
			\item Se apenas I e III estiverem corretas.
			\item Se apenas II e IV estiverem corretas.
			\item \alert<2>{Se apenas I, II e III estiverem corretas.}
			\item Se apenas IV estiver correta.
			\item Se todas estiverem corretas.
		\end{enumerate}
	\end{enumerate}
\end{frame}

\begin{frame}{Princípios do Gerenciamento de Projetos}
	\begin{enumerate}
		\setcounter{enumi}{9}
		\justifying
		\item Relacione os elementos da colaboração às suas características:
		\begin{columns}
			\begin{column}{0.32\textwidth}
				\begin{enumerate}[A]
					\item Acordos de Equipe
					\item Estruturas Organizacionais
					\item Processos da Equipe
					\item Diversidade
				\end{enumerate}
			\end{column}
			\begin{column}{0.68\textwidth}
				\begin{itemize}
					\item Conjunto de normas e diretrizes comportamentais criadas pela própria equipe.
					\item A forma como a empresa está montada que impacta a autonomia.
					\item Como o trabalho é realizado e como as decisões são tomadas.
					\item Diferentes perspectivas e habilidades que enriquecem a solução.
				\end{itemize}
			\end{column}
		\end{columns}
		\begin{enumerate}[a]
			\item A, B, C, D.
			\item \alert<2>{B, D, C, A}
			\item C, A, D, B.
			\item D, C, B, A.
			\item B, C, D, A.
		\end{enumerate}
	\end{enumerate}
\end{frame}

\begin{frame}{Princípios do Gerenciamento de Projetos}
	\begin{enumerate}
		\setcounter{enumi}{10}
		\justifying
		\item Em ambientes colaborativos, o líder deve atuar principalmente como:
		\begin{enumerate}[a]
			\item Um juiz autoritário.
			\item \alert<2>{Um facilitador que remove impedimentos.}
			\item Um controlador de microtarefas.
			\item Um porta-voz que isola a equipe.
			\item Um técnico que refaz o trabalho alheio.
		\end{enumerate}
	\end{enumerate}
\end{frame}

\begin{frame}{Princípios do Gerenciamento de Projetos}
	\begin{enumerate}
		\setcounter{enumi}{11}
		\justifying
		\item Um membro júnior sugere uma inovação, mas um sênior o interrompe rudemente. Para manter a colaboração, você deve:
		\begin{enumerate}[a]
			\item Concordar com o sênior.
			\item Ignorar o comentário.
			\item \alert<2>{Facilitar um diálogo focado no mérito da ideia e no respeito mútuo.}
			\item Punir o júnior.
			\item Decidir sozinho para encerrar a briga.
		\end{enumerate}
	\end{enumerate}
\end{frame}

\begin{frame}{Princípios do Gerenciamento de Projetos}
	\begin{enumerate}
		\setcounter{enumi}{12}
		\justifying
		\item Sobre as responsabilidades da equipe colaborativa:
		\begin{enumerate}[I]
			\item Responsabilidade individual apenas por tarefas próprias.
			\item Responsabilidade compartilhada pelos resultados.
			\item Estímulo ao cross-training (ajuda mútua).
			\item Comunicação transparente e honesta.
		\end{enumerate}
		\begin{enumerate}[a]
			\item Se apenas II, III e IV estiverem corretas.
			\item Se apenas I e III estiverem corretas.
			\item \alert<2>{Se apenas II, III e IV estiverem corretas.}
			\item Se apenas IV estiver correta.
			\item Se todas estiverem corretas.
		\end{enumerate}
	\end{enumerate}
\end{frame}

\begin{frame}{Princípios do Gerenciamento de Projetos}
	\begin{enumerate}
		\setcounter{enumi}{13}
		\justifying
		\item Qual influência das partes interessadas é crítica para o alcance do valor?
		\begin{enumerate}[a]
			\item Registro de horas trabalhadas.
			\item \alert<2>{Definição e aceitação dos requisitos e critérios de sucesso.}
			\item Escolha da marca do café no escritório.
			\item Preenchimento de formulários de férias.
			\item Manutenção de arquivos físicos redundantes.
		\end{enumerate}
	\end{enumerate}
\end{frame}

\begin{frame}{Princípios do Gerenciamento de Projetos}
	\begin{enumerate}
		\setcounter{enumi}{14}
		\justifying
		\item O envolvimento exige habilidades como a \rule{60pt}{1pt} ativa e a gestão de \rule{60pt}{1pt}.
		\begin{enumerate}[a]
			\item Comunicação / Tempo
			\item \alert<2>{Escuta / Conflitos}
			\item Negociação / Riscos
			\item Apresentação / Poder
			\item Leitura / Documentação
		\end{enumerate}
	\end{enumerate}
\end{frame}

\begin{frame}{Princípios do Gerenciamento de Projetos}
	\begin{enumerate}
		\setcounter{enumi}{15}
		\justifying
		\item Atividades de engajamento de stakeholders incluem:
		\begin{enumerate}[I]
			\item Identificação de interesses.
			\item Análise de poder e influência.
			\item Monitoramento de mudanças de cenário.
			\item Definição de estratégias de comunicação.
		\end{enumerate}
		\begin{enumerate}[a]
			\item Se todas estiverem corretas.
			\item Se apenas I e III estiverem corretas.
			\item \alert<2>{Se todas estiverem corretas.}
			\item Se apenas I, II e IV estiverem corretas.
			\item Se apenas IV estiver correta.
		\end{enumerate}
	\end{enumerate}
\end{frame}

\begin{frame}{Princípios do Gerenciamento de Projetos}
	\begin{enumerate}
		\setcounter{enumi}{16}
		\justifying
		\item Relacione a ação ao objetivo:
		\begin{columns}
			\begin{column}{0.28\textwidth}
				\begin{enumerate}[A]
					\item Identificação
					\item Análise
					\item Monitoramento
					\item Engajamento
				\end{enumerate}
			\end{column}
			\begin{column}{0.72\textwidth}
				\begin{itemize}
					\item Descobrir quem são os afetados.
					\item Entender as motivações e o apoio.
					\item Verificar mudanças de interesse ou poder.
					\item Interagir para resolver problemas.
				\end{itemize}
			\end{column}
		\end{columns}
		\begin{enumerate}[a]
			\item A, B, C, D.
			\item \alert<2>{B, A, D, C}
			\item C, D, A, B.
			\item D, C, B, A.
			\item A, C, B, D.
		\end{enumerate}
	\end{enumerate}
\end{frame}

\begin{frame}{Princípios do Gerenciamento de Projetos}
	\begin{enumerate}
		\setcounter{enumi}{17}
		\justifying
		\item Comunicação bidirecional significa:
		\begin{enumerate}[a]
			\item Enviar e-mail duas vezes.
			\item \alert<2>{Transmissão de informação + escuta/feedback.}
			\item Somente o patrocinador fala.
			\item Usar dois idiomas.
			\item Usar apenas canais digitais.
		\end{enumerate}
	\end{enumerate}
\end{frame}

\begin{frame}{Princípios do Gerenciamento de Projetos}
	\begin{enumerate}
		\setcounter{enumi}{18}
		\justifying
		\item Reuniões públicas para ouvir a comunidade sobre ruídos de uma obra é exemplo de:
		\begin{enumerate}[a]
			\item Gestão de Escopo.
			\item \alert<2>{Engajamento proativo de stakeholders.}
			\item Adiar o projeto.
			\item Transferência de risco.
			\item Microgerenciamento.
		\end{enumerate}
	\end{enumerate}
\end{frame}

\begin{frame}{Princípios do Gerenciamento de Projetos}
	\begin{enumerate}
		\setcounter{enumi}{19}
		\justifying
		\item São exemplos de partes interessadas internas:
		\begin{enumerate}[I]
			\item Equipe do projeto.
			\item Patrocinador.
			\item Órgãos reguladores governamentais.
			\item Gerentes funcionais.
		\end{enumerate}
		\begin{enumerate}[a]
			\item Se apenas I, II e IV estiverem corretas.
			\item Se apenas II e III estiverem corretas.
			\item \alert<2>{Se apenas I, II e IV estiverem corretas.}
			\item Se apenas IV estiver correta.
			\item Se todas estiverem corretas.
		\end{enumerate}
	\end{enumerate}
\end{frame}

\begin{frame}{Princípios do Gerenciamento de Projetos}
	\begin{enumerate}
		\setcounter{enumi}{20}
		\justifying
		\item Se o valor previsto não for mais alcançável, a ação é:
		\begin{enumerate}[a]
			\item Continuar o projeto.
			\item Ignorar o fato.
			\item \alert<2>{Recomendar alteração ou encerramento.}
			\item Esconder a informação.
			\item Dobrar o escopo.
		\end{enumerate}
	\end{enumerate}
\end{frame}

\begin{frame}{Princípios do Gerenciamento de Projetos}
	\begin{enumerate}
		\setcounter{enumi}{21}
		\justifying
		\item A entrega de valor envolve a transição do estado atual para um \rule{60pt}{1pt} \rule{60pt}{1pt} desejado.
		\begin{enumerate}[a]
			\item Estado / Atual
			\item \alert<2>{Estado Futuro}
			\item Resultado / Final
			\item Benefício / Esperado
			\item Output / Tangível
		\end{enumerate}
	\end{enumerate}
\end{frame}

\begin{frame}{Princípios do Gerenciamento de Projetos}
	\begin{enumerate}
		\setcounter{enumi}{22}
		\justifying
		\item Contribuem para o foco no valor:
		\begin{enumerate}[I]
			\item Alinhamento estratégico.
			\item Identificação de benefícios.
			\item Priorização apenas técnica.
			\item Avaliação da viabilidade econômica.
		\end{enumerate}
		\begin{enumerate}[a]
			\item Se apenas I, II e IV estiverem corretas.
			\item Se apenas I e III estiverem corretas.
			\item \alert<2>{Se apenas I, II e IV estiverem corretas.}
			\item Se apenas IV estiver correta.
			\item Se todas estiverem corretas.
		\end{enumerate}
	\end{enumerate}
\end{frame}

\begin{frame}{Princípios do Gerenciamento de Projetos}
	\begin{enumerate}
		\setcounter{enumi}{23}
		\justifying
		\item Relacione os conceitos:
		\begin{columns}
			\begin{column}{0.32\textwidth}
				\begin{enumerate}[A]
					\item Valor de Negócio
					\item Benefício
					\item Resultado (Outcome)
					\item Output (Saída)
				\end{enumerate}
			\end{column}
			\begin{column}{0.68\textwidth}
				\begin{itemize}
					\item Utilidade ou mérito para a organização.
					\item Ganho obtido com o resultado.
					\item Efeito ou consequência do projeto.
					\item Produto ou serviço tangível criado.
				\end{itemize}
			\end{column}
		\end{columns}
		\begin{enumerate}[a]
			\item A, B, C, D.
			\item \alert<2>{A, B, C, D}
			\item C, D, A, B.
			\item D, C, B, A.
			\item B, A, D, C.
		\end{enumerate}
	\end{enumerate}
\end{frame}

\begin{frame}{Princípios do Gerenciamento de Projetos}
	\begin{enumerate}
		\setcounter{enumi}{24}
		\justifying
		\item Sistema de segurança cibernética gera valor quando:
		\begin{enumerate}[a]
			\item Código é sem erros.
			\item Manual é impresso.
			\item \alert<2>{Há menos ataques e custos reduzidos.}
			\item Gerente ganha certificação.
			\item Equipe faz horas extras.
		\end{enumerate}
	\end{enumerate}
\end{frame}

\begin{frame}{Princípios do Gerenciamento de Projetos}
	\begin{enumerate}
		\setcounter{enumi}{25}
		\justifying
		\item Matéria-prima dobra de preço e ROI fica negativo. O gerente deve:
		\begin{enumerate}[a]
			\item Trabalhar mais rápido.
			\item \alert<2>{Solicitar reunião de governança para reavaliar viabilidade.}
			\item Alterar dados financeiros.
			\item Culpar compras.
			\item Esperar o projeto acabar.
		\end{enumerate}
	\end{enumerate}
\end{frame}

\begin{frame}{Princípios do Gerenciamento de Projetos}
	\begin{enumerate}
		\setcounter{enumi}{26}
		\justifying
		\item A qualidade no PMBOK 7 é entendida como:
		\begin{enumerate}[a]
			\item Atender apenas aos requisitos técnicos do produto.
			\item \alert<2>{Atender aos requisitos e às expectativas das partes interessadas.}
			\item Entregar o mais rápido possível.
			\item Reduzir ao máximo os custos.
			\item Seguir estritamente o plano original.
		\end{enumerate}
	\end{enumerate}
\end{frame}

\begin{frame}{Princípios do Gerenciamento de Projetos}
	\begin{enumerate}
		\setcounter{enumi}{27}
		\justifying
		\item O custo da qualidade inclui:
		\begin{enumerate}[I]
			\item Custos de prevenção.
			\item Custos de avaliação (inspeção e testes).
			\item Custos de falhas internas.
			\item Custos de falhas externas.
		\end{enumerate}
		\begin{enumerate}[a]
			\item Se apenas I e II estiverem corretas.
			\item \alert<2>{Se todas estiverem corretas.}
			\item Se apenas III e IV estiverem corretas.
			\item Se apenas I, II e III estiverem corretas.
			\item Se apenas II estiver correta.
		\end{enumerate}
	\end{enumerate}
\end{frame}

\begin{frame}{Princípios do Gerenciamento de Projetos}
	\begin{enumerate}
		\setcounter{enumi}{28}
		\justifying
		\item Relacione os conceitos de qualidade:
		\begin{columns}
			\begin{column}{0.35\textwidth}
				\begin{enumerate}[A]
					\item Prevenção
					\item Avaliação
					\item Falha Interna
					\item Falha Externa
				\end{enumerate}
			\end{column}
			\begin{column}{0.65\textwidth}
				\begin{itemize}
					\item Treinamento, processos robustos e planejamento.
					\item Inspeções, testes e revisões.
					\item Defeitos encontrados antes da entrega ao cliente.
					\item Defeitos encontrados após a entrega ao cliente.
				\end{itemize}
			\end{column}
		\end{columns}
		\begin{enumerate}[a]
			\item A, B, C, D.
			\item \alert<2>{A, B, C, D}
			\item C, D, A, B.
			\item D, C, B, A.
			\item B, A, D, C.
		\end{enumerate}
	\end{enumerate}
\end{frame}

\begin{frame}{Princípios do Gerenciamento de Projetos}
	\begin{enumerate}
		\setcounter{enumi}{29}
		\justifying
		\item "Fazer certo da primeira vez" alinha-se com:
		\begin{enumerate}[a]
			\item Inspeção em massa.
			\item Custos variáveis.
			\item \alert<2>{Prevenção de defeitos.}
			\item Aceitação de falhas.
			\item Eliminação de escopo.
		\end{enumerate}
	\end{enumerate}
\end{frame}

\begin{frame}{Princípios do Gerenciamento de Projetos}
	\begin{enumerate}
		\setcounter{enumi}{30}
		\justifying
		\item Pular testes para recuperar atraso resultou em erros no cliente. Esta decisão foi:
		\begin{enumerate}[a]
			\item Eficaz pelo cronograma.
			\item \alert<2>{Incorreta, custo do defeito pós-entrega é superior à prevenção.}
			\item Irrelevante.
			\item Boa adaptação.
			\item Culpa do cliente.
		\end{enumerate}
	\end{enumerate}
\end{frame}

\begin{frame}{Princípios do Gerenciamento de Projetos}
	\begin{enumerate}
		\setcounter{enumi}{31}
		\justifying
		\item Sobre Processos de Qualidade:
		\begin{enumerate}[I]
			\item Adaptados à complexidade.
			\item Documentados para melhoria.
			\item Incluem revisões e testes.
			\item Ignorados se houver pressa.
		\end{enumerate}
		\begin{enumerate}[a]
			\item Se apenas I, II e III estiverem corretas.
			\item Se apenas II e IV estiverem corretas.
			\item \alert<2>{Se apenas I, II e III estiverem corretas.}
			\item Se apenas IV estiver correta.
			\item Se todas estiverem corretas.
		\end{enumerate}
	\end{enumerate}
\end{frame}

\begin{frame}{Princípios do Gerenciamento de Projetos}
	\begin{enumerate}
		\setcounter{enumi}{32}
		\justifying
		\item As três fontes primárias de complexidade são: o comportamento \rule{50pt}{1pt}, o comportamento do \rule{50pt}{1pt} e a \rule{50pt}{1pt}.
		\begin{enumerate}[a]
			\item Humano, Sistema, Ambiguidade
			\item \alert<2>{Humano, Sistema, Ambiguidade}
			\item Técnico, Organizacional, Externo
			\item Financeiro, Temporal, Escopo
			\item Social, Cultural, Político
		\end{enumerate}
	\end{enumerate}
\end{frame}

\begin{frame}{Princípios do Gerenciamento de Projetos}
	\begin{enumerate}
		\setcounter{enumi}{33}
		\justifying
		\item Quais fatores contribuem para a Complexidade de Sistema?
		\begin{enumerate}[I]
			\item Interdependências entre componentes internos e externos.
			\item Mudanças frequentes nos requisitos.
			\item Desconexão entre estratégia e execução.
			\item Uso de uma única folha de cálculo para controlo.
		\end{enumerate}
		\begin{enumerate}[a]
			\item Se apenas I, II e III estiverem corretas.
			\item Se apenas II e IV estiverem corretas.
			\item \alert<2>{Se apenas I, II e III estiverem corretas.}
			\item Se apenas IV estiver correta.
			\item Se todas estiverem corretas.
		\end{enumerate}
	\end{enumerate}
\end{frame}

\begin{frame}{Princípios do Gerenciamento de Projetos}
	\begin{enumerate}
		\setcounter{enumi}{34}
		\justifying
		\item Relacione as fontes de complexidade aos exemplos:
		\begin{columns}
			\begin{column}{0.32\textwidth}
				\begin{enumerate}[A]
					\item Comportamento Humano
					\item Ambiguidade
					\item Complexidade Tecnológica
					\item Incerteza
				\end{enumerate}
			\end{column}
			\begin{column}{0.68\textwidth}
				\begin{itemize}
					\item Conflitos de interesses e dinâmicas de poder.
					\item Falta de clareza sobre o significado de um evento.
					\item Integração de IA nunca antes testada.
					\item Probabilidade de eventos futuros desconhecidos.
				\end{itemize}
			\end{column}
		\end{columns}
		\begin{enumerate}[a]
			\item A, B, C, D.
			\item \alert<2>{A, B, C, D}
			\item C, D, A, B.
			\item D, C, B, A.
			\item B, A, D, C.
		\end{enumerate}
	\end{enumerate}
\end{frame}

\begin{frame}{Princípios do Gerenciamento de Projetos}
	\begin{enumerate}
		\setcounter{enumi}{35}
		\justifying
		\item Ações ao encontrar alta complexidade:
		\begin{enumerate}[I]
			\item Identificar os elementos geradores de complexidade.
			\item Analisar as interconexões e impactos sistémicos.
			\item Adaptar a abordagem de gestão.
			\item Monitorar e reavaliar continuamente.
		\end{enumerate}
		\begin{enumerate}[a]
			\item 4, 2, 1, 3.
			\item \alert<2>{1, 2, 3, 4}
			\item 2, 1, 4, 3.
			\item 3, 4, 1, 2.
			\item 1, 3, 2, 4.
		\end{enumerate}
	\end{enumerate}
\end{frame}

\begin{frame}{Princípios do Gerenciamento de Projetos}
	\begin{enumerate}
		\setcounter{enumi}{36}
		\justifying
		\item Para navegar na complexidade, as equipes devem focar em:
		\begin{enumerate}[a]
			\item Controlo total de variáveis por contratos rígidos.
			\item Eliminar requisitos difíceis.
			\item \alert<2>{Desenvolver resiliência e aprendizagem rápida.}
			\item Ignorar a complexidade.
			\item Delegar tudo a consultores externos.
		\end{enumerate}
	\end{enumerate}
\end{frame}

\begin{frame}{Princípios do Gerenciamento de Projetos}
	\begin{enumerate}
		\setcounter{enumi}{37}
		\justifying
		\item Fusão de empresas com resistência cultural severa é complexidade de:
		\begin{enumerate}[a]
			\item Escopo.
			\item \alert<2>{Comportamento Humano.}
			\item Orçamento.
			\item Hardware.
			\item Cronograma.
		\end{enumerate}
	\end{enumerate}
\end{frame}

\begin{frame}{Princípios do Gerenciamento de Projetos}
	\begin{enumerate}
		\setcounter{enumi}{38}
		\justifying
		\item Estratégias para mitigar efeitos da complexidade:
		\begin{enumerate}[I]
			\item Descentralização da tomada de decisão.
			\item Comunicação transparente e frequente.
			\item Experimentos e prototipagem rápida.
			\item Abordagem estritamente preditiva em alta incerteza.
		\end{enumerate}
		\begin{enumerate}[a]
			\item Se apenas I, II e III estiverem corretas.
			\item Se apenas II e IV estiverem corretas.
			\item \alert<2>{Se apenas I, II e III estiverem corretas.}
			\item Se apenas IV estiver correta.
			\item Se apenas I e III estiverem corretas.
		\end{enumerate}
	\end{enumerate}
\end{frame}

\begin{frame}{Princípios do Gerenciamento de Projetos}
	\begin{enumerate}
		\setcounter{enumi}{39}
		\justifying
		\item Para ameaças: evitar, transferir, \rule{50pt}{1pt}, aceitar. Para oportunidades: explorar, compartilhar, \rule{50pt}{1pt}, aceitar.
		\begin{enumerate}[a]
			\item Mitigar / Melhorar (Enhance)
			\item \alert<2>{Mitigar / Melhorar (Enhance)}
			\item Aceitar / Explorar
			\item Transferir / Compartilhar
			\item Evitar / Melhorar
		\end{enumerate}
	\end{enumerate}
\end{frame}

\begin{frame}{Princípios do Gerenciamento de Projetos}
	\begin{enumerate}
		\setcounter{enumi}{40}
		\justifying
		\item Fatores na escolha da estratégia de risco:
		\begin{enumerate}[I]
			\item Probabilidade.
			\item Impacto.
			\item Custo-benefício da resposta.
			\item Cor favorita do sponsor.
		\end{enumerate}
		\begin{enumerate}[a]
			\item Se apenas I, II e III estiverem corretas.
			\item Se apenas II e IV estiverem corretas.
			\item \alert<2>{Se apenas I, II e III estiverem corretas.}
			\item Se apenas IV estiver correta.
			\item Se apenas I e III estiverem corretas.
		\end{enumerate}
	\end{enumerate}
\end{frame}

\begin{frame}{Princípios do Gerenciamento de Projetos}
	\begin{enumerate}
		\setcounter{enumi}{41}
		\justifying
		\item Relacione estratégias de ameaça:
		\begin{columns}
			\begin{column}{0.28\textwidth}
				\begin{enumerate}[A]
					\item Evitar
					\item Transferir
					\item Mitigar
					\item Aceitar
				\end{enumerate}
			\end{column}
			\begin{column}{0.72\textwidth}
				\begin{itemize}
					\item Eliminar atividade perigosa.
					\item Contratar seguro.
					\item Testes extras para reduzir falha.
					\item Reserva de contingência.
				\end{itemize}
			\end{column}
		\end{columns}
		\begin{enumerate}[a]
			\item A, B, C, D.
			\item \alert<2>{B, A, D, C}
			\item C, D, A, B.
			\item D, C, B, A.
			\item A, C, B, D.
		\end{enumerate}
	\end{enumerate}
\end{frame}

\begin{frame}{Princípios do Gerenciamento de Projetos}
	\begin{enumerate}
		\setcounter{enumi}{42}
		\justifying
		\item Etapas da otimização de riscos:
		\begin{enumerate}[I]
			\item Identificar riscos.
			\item Análise qualitativa/quantitativa.
			\item Planear e implementar respostas.
			\item Monitorar riscos residuais.
		\end{enumerate}
		\begin{enumerate}[a]
			\item 4, 2, 1, 3.
			\item \alert<2>{1, 2, 3, 4}
			\item 2, 1, 4, 3.
			\item 3, 4, 1, 2.
			\item 1, 3, 2, 4.
		\end{enumerate}
	\end{enumerate}
\end{frame}

\begin{frame}{Princípios do Gerenciamento de Projetos}
	\begin{enumerate}
		\setcounter{enumi}{43}
		\justifying
		\item Explorar (oportunidades) visa:
		\begin{enumerate}[a]
			\item Dividir responsabilidade.
			\item Aumentar probabilidade/impacto.
			\item \alert<2>{Garantir que a oportunidade ocorra.}
			\item Ignorar.
			\item Atribuir ao jurídico.
		\end{enumerate}
	\end{enumerate}
\end{frame}

\begin{frame}{Princípios do Gerenciamento de Projetos}
	\begin{enumerate}
		\setcounter{enumi}{44}
		\justifying
		\item Antecipar compra com dólar baixo para garantir preço é:
		\begin{enumerate}[a]
			\item Ameaça mitigada.
			\item Incerteza evitada.
			\item \alert<2>{Oportunidade explorada.}
			\item Mudança aceite.
			\item Falha de governança.
		\end{enumerate}
	\end{enumerate}
\end{frame}

\begin{frame}{Princípios do Gerenciamento de Projetos}
	\begin{enumerate}
		\setcounter{enumi}{45}
		\justifying
		\item Sobre Riscos Residuais:
		\begin{enumerate}[I]
			\item Permanecem após a resposta.
			\item Devem ser monitorados.
			\item Devem ser zerados sempre.
			\item Podem exigir planos de contingência.
		\end{enumerate}
		\begin{enumerate}[a]
			\item Se apenas I, II e IV estiverem corretas.
			\item Se apenas II e III estiverem corretas.
			\item \alert<2>{Se apenas I, II e IV estiverem corretas.}
			\item Se apenas IV estiver correta.
			\item Se apenas I e III estiverem corretas.
		\end{enumerate}
	\end{enumerate}
\end{frame}

\begin{frame}{Princípios do Gerenciamento de Projetos}
	\begin{enumerate}
		\setcounter{enumi}{46}
		\justifying
		\item A \rule{50pt}{0.75pt} permite ajustar-se a novas situações; a \rule{50pt}{0.75pt} garante voltar ao trilho após falhas.
		\begin{enumerate}[a]
			\item Resiliência / Adaptabilidade
			\item \alert<2>{Adaptabilidade / Resiliência}
			\item Flexibilidade / Rigidez
			\item Mudança / Estabilidade
			\item Velocidade / Precisão
		\end{enumerate}
	\end{enumerate}
\end{frame}

\begin{frame}{Princípios do Gerenciamento de Projetos}
	\begin{enumerate}
		\setcounter{enumi}{47}
		\justifying
		\item Favorecem adaptabilidade e resiliência:
		\begin{enumerate}[I]
			\item Transparência.
			\item Cultura de experimentação.
			\item Centralização total.
			\item Soluções modulares.
		\end{enumerate}
		\begin{enumerate}[a]
			\item Se apenas I, II e IV estiverem corretas.
			\item Se apenas II e IV estiverem corretas.
			\item \alert<2>{Se apenas I, II e IV estiverem corretas.}
			\item Se apenas IV estiver correta.
			\item Se apenas I e III estiverem corretas.
		\end{enumerate}
	\end{enumerate}
\end{frame}

\begin{frame}{Princípios do Gerenciamento de Projetos}
	\begin{enumerate}
		\setcounter{enumi}{48}
		\justifying
		\item Relacione conceitos à prática:
		\begin{columns}
			\begin{column}{0.3\textwidth}
				\begin{enumerate}[A]
					\item Adaptabilidade
					\item Resiliência
					\item Feedback
					\item Inspeção
				\end{enumerate}
			\end{column}
			\begin{column}{0.7\textwidth}
				\begin{itemize}
					\item Ajustar design após teste.
					\item Recuperar de atraso por desastre natural.
					\item Informação para validar direção.
					\item Verificação técnica de padrões.
				\end{itemize}
			\end{column}
		\end{columns}
		\begin{enumerate}[a]
			\item A, B, C, D.
			\item \alert<2>{B, A, C, D}
			\item C, D, A, B.
			\item D, C, B, A.
			\item A, C, B, D.
		\end{enumerate}
	\end{enumerate}
\end{frame}

\begin{frame}{Princípios do Gerenciamento de Projetos}
	\begin{enumerate}
		\setcounter{enumi}{49}
		\justifying
		\item Desenvolver resiliência perante problema:
		\begin{enumerate}[I]
			\item Aceitar a ocorrência.
			\item Estabilizar e avaliar danos.
			\item Encontrar caminhos alternativos.
			\item Extrair lições aprendidas.
		\end{enumerate}
		\begin{enumerate}[a]
			\item 4, 2, 1, 3.
			\item \alert<2>{1, 2, 3, 4}
			\item 2, 1, 4, 3.
			\item 3, 4, 1, 2.
			\item 1, 3, 2, 4.
		\end{enumerate}
	\end{enumerate}
\end{frame}

\begin{frame}{Princípios do Gerenciamento de Projetos}
	\begin{enumerate}
		\setcounter{enumi}{50}
		\justifying
		\item Abordagem adaptativa abraça mudança porque:
		\begin{enumerate}[a]
			\item Não exige planeamento.
			\item Custo é irrelevante.
			\item \alert<2>{Requisitos totais podem ser desconhecidos.}
			\item Projeto nunca termina.
			\item Não exige documentação.
		\end{enumerate}
	\end{enumerate}
\end{frame}

\begin{frame}{Princípios do Gerenciamento de Projetos}
	\begin{enumerate}
		\setcounter{enumi}{51}
		\justifying
		\item Equipe pivota projeto após lançamento de concorrente mais barato. Isso é:
		\begin{enumerate}[a]
			\item Falta de planeamento.
			\item \alert<2>{Alta adaptabilidade.}
			\item Desperdício.
			\item Incompetência.
			\item Falha de governança.
		\end{enumerate}
	\end{enumerate}
\end{frame}

\begin{frame}{Princípios do Gerenciamento de Projetos}
	\begin{enumerate}
		\setcounter{enumi}{52}
		\justifying
		\item Dificultam a resiliência:
		\begin{enumerate}[I]
			\item Punição a erros honestos.
			\item Burocracia lenta.
			\item Silos de informação.
			\item Autonomia delegada.
		\end{enumerate}
		\begin{enumerate}[a]
			\item Se apenas I, II e III estiverem corretas.
			\item Se apenas II e IV estiverem corretas.
			\item \alert<2>{Se apenas I, II e III estiverem corretas.}
			\item Se apenas IV estiver correta.
			\item Se apenas I e III estiverem corretas.
		\end{enumerate}
	\end{enumerate}
\end{frame}

\begin{frame}{Princípios do Gerenciamento de Projetos}
	\begin{enumerate}
		\setcounter{enumi}{53}
		\justifying
		\item As três fontes primárias de complexidade são: o comportamento \rule{50pt}{1pt}, o comportamento do \rule{50pt}{1pt} e a \rule{50pt}{1pt}.
		\begin{enumerate}[a]
			\item Humano, Sistema, Ambiguidade
			\item \alert<2>{Humano, Sistema, Ambiguidade}
			\item Técnico, Organizacional, Externo
			\item Financeiro, Temporal, Escopo
			\item Social, Cultural, Político
		\end{enumerate}
	\end{enumerate}
\end{frame}

\begin{frame}{Princípios do Gerenciamento de Projetos}
	\begin{enumerate}
		\setcounter{enumi}{54}
		\justifying
		\item Quais fatores contribuem para a Complexidade de Sistema?
		\begin{enumerate}[I]
			\item Interdependências entre componentes internos e externos.
			\item Mudanças frequentes nos requisitos.
			\item Desconexão entre estratégia e execução.
			\item Uso de uma única folha de cálculo para controlo.
		\end{enumerate}
		\begin{enumerate}[a]
			\item Se apenas I, II e III estiverem corretas.
			\item Se apenas II e IV estiverem corretas.
			\item \alert<2>{Se apenas I, II e III estiverem corretas.}
			\item Se apenas IV estiver correta.
			\item Se todas estiverem corretas.
		\end{enumerate}
	\end{enumerate}
\end{frame}

\begin{frame}{Princípios do Gerenciamento de Projetos}
	\begin{enumerate}
		\setcounter{enumi}{55}
		\justifying
		\item Relacione as fontes de complexidade aos exemplos:
		\begin{columns}
			\begin{column}{0.32\textwidth}
				\begin{enumerate}[A]
					\item Comportamento Humano
					\item Ambiguidade
					\item Complexidade Tecnológica
					\item Incerteza
				\end{enumerate}
			\end{column}
			\begin{column}{0.68\textwidth}
				\begin{itemize}
					\item Conflitos de interesses e dinâmicas de poder.
					\item Falta de clareza sobre o significado de um evento.
					\item Integração de IA nunca antes testada.
					\item Probabilidade de eventos futuros desconhecidos.
				\end{itemize}
			\end{column}
		\end{columns}
		\begin{enumerate}[a]
			\item A, B, C, D.
			\item \alert<2>{A, B, C, D}
			\item C, D, A, B.
			\item D, C, B, A.
			\item B, A, D, C.
		\end{enumerate}
	\end{enumerate}
\end{frame}

\begin{frame}{Princípios do Gerenciamento de Projetos}
	\begin{enumerate}
		\setcounter{enumi}{56}
		\justifying
		\item Ações ao encontrar alta complexidade:
		\begin{enumerate}[I]
			\item Identificar os elementos geradores de complexidade.
			\item Analisar as interconexões e impactos sistémicos.
			\item Adaptar a abordagem de gestão.
			\item Monitorar e reavaliar continuamente.
		\end{enumerate}
		\begin{enumerate}[a]
			\item 4, 2, 1, 3.
			\item \alert<2>{1, 2, 3, 4}
			\item 2, 1, 4, 3.
			\item 3, 4, 1, 2.
			\item 1, 3, 2, 4.
		\end{enumerate}
	\end{enumerate}
\end{frame}

\begin{frame}{Princípios do Gerenciamento de Projetos}
	\begin{enumerate}
		\setcounter{enumi}{57}
		\justifying
		\item Para navegar na complexidade, as equipes devem focar em:
		\begin{enumerate}[a]
			\item Controlo total de variáveis por contratos rígidos.
			\item Eliminar requisitos difíceis.
			\item \alert<2>{Desenvolver resiliência e aprendizagem rápida.}
			\item Ignorar a complexidade.
			\item Delegar tudo a consultores externos.
		\end{enumerate}
	\end{enumerate}
\end{frame}

\begin{frame}{Princípios do Gerenciamento de Projetos}
	\begin{enumerate}
		\setcounter{enumi}{58}
		\justifying
		\item Fusão de empresas com resistência cultural severa é complexidade de:
		\begin{enumerate}[a]
			\item Escopo.
			\item \alert<2>{Comportamento Humano.}
			\item Orçamento.
			\item Hardware.
			\item Cronograma.
		\end{enumerate}
	\end{enumerate}
\end{frame}

\begin{frame}{Princípios do Gerenciamento de Projetos}
	\begin{enumerate}
		\setcounter{enumi}{59}
		\justifying
		\item Estratégias para mitigar efeitos da complexidade:
		\begin{enumerate}[I]
			\item Descentralização da tomada de decisão.
			\item Comunicação transparente e frequente.
			\item Experimentos e prototipagem rápida.
			\item Abordagem estritamente preditiva em alta incerteza.
		\end{enumerate}
		\begin{enumerate}[a]
			\item Se apenas I, II e III estiverem corretas.
			\item Se apenas II e IV estiverem corretas.
			\item \alert<2>{Se apenas I, II e III estiverem corretas.}
			\item Se apenas IV estiver correta.
			\item Se apenas I e III estiverem corretas.
		\end{enumerate}
	\end{enumerate}
\end{frame}

\begin{frame}{Princípios do Gerenciamento de Projetos}
	\begin{enumerate}
		\setcounter{enumi}{60}
		\justifying
		\item Para ameaças: evitar, transferir, \rule{50pt}{1pt}, aceitar. Para oportunidades: explorar, compartilhar, \rule{50pt}{1pt}, aceitar.
		\begin{enumerate}[a]
			\item Mitigar / Melhorar (Enhance)
			\item \alert<2>{Mitigar / Melhorar (Enhance)}
			\item Aceitar / Explorar
			\item Transferir / Compartilhar
			\item Evitar / Melhorar
		\end{enumerate}
	\end{enumerate}
\end{frame}

\begin{frame}{Princípios do Gerenciamento de Projetos}
	\begin{enumerate}
		\setcounter{enumi}{61}
		\justifying
		\item Fatores na escolha da estratégia de risco:
		\begin{enumerate}[I]
			\item Probabilidade.
			\item Impacto.
			\item Custo-benefício da resposta.
			\item Cor favorita do sponsor.
		\end{enumerate}
		\begin{enumerate}[a]
			\item Se apenas I, II e III estiverem corretas.
			\item Se apenas II e IV estiverem corretas.
			\item \alert<2>{Se apenas I, II e III estiverem corretas.}
			\item Se apenas IV estiver correta.
			\item Se apenas I e III estiverem corretas.
		\end{enumerate}
	\end{enumerate}
\end{frame}

\begin{frame}{Princípios do Gerenciamento de Projetos}
	\begin{enumerate}
		\setcounter{enumi}{62}
		\justifying
		\item Relacione estratégias de ameaça:
		\begin{columns}
			\begin{column}{0.28\textwidth}
				\begin{enumerate}[A]
					\item Evitar
					\item Transferir
					\item Mitigar
					\item Aceitar
				\end{enumerate}
			\end{column}
			\begin{column}{0.72\textwidth}
				\begin{itemize}
					\item Eliminar atividade perigosa.
					\item Contratar seguro.
					\item Testes extras para reduzir falha.
					\item Reserva de contingência.
				\end{itemize}
			\end{column}
		\end{columns}
		\begin{enumerate}[a]
			\item A, B, C, D.
			\item \alert<2>{B, A, D, C}
			\item C, D, A, B.
			\item D, C, B, A.
			\item A, C, B, D.
		\end{enumerate}
	\end{enumerate}
\end{frame}

\begin{frame}{Princípios do Gerenciamento de Projetos}
	\begin{enumerate}
		\setcounter{enumi}{63}
		\justifying
		\item Etapas da otimização de riscos:
		\begin{enumerate}[I]
			\item Identificar riscos.
			\item Análise qualitativa/quantitativa.
			\item Planear e implementar respostas.
			\item Monitorar riscos residuais.
		\end{enumerate}
		\begin{enumerate}[a]
			\item 4, 2, 1, 3.
			\item \alert<2>{1, 2, 3, 4}
			\item 2, 1, 4, 3.
			\item 3, 4, 1, 2.
			\item 1, 3, 2, 4.
		\end{enumerate}
	\end{enumerate}
\end{frame}

\begin{frame}{Princípios do Gerenciamento de Projetos}
	\begin{enumerate}
		\setcounter{enumi}{64}
		\justifying
		\item Explorar (oportunidades) visa:
		\begin{enumerate}[a]
			\item Dividir responsabilidade.
			\item Aumentar probabilidade/impacto.
			\item \alert<2>{Garantir que a oportunidade ocorra.}
			\item Ignorar.
			\item Atribuir ao jurídico.
		\end{enumerate}
	\end{enumerate}
\end{frame}

\begin{frame}{Princípios do Gerenciamento de Projetos}
	\begin{enumerate}
		\setcounter{enumi}{65}
		\justifying
		\item Antecipar compra com dólar baixo para garantir preço é:
		\begin{enumerate}[a]
			\item Ameaça mitigada.
			\item Incerteza evitada.
			\item \alert<2>{Oportunidade explorada.}
			\item Mudança aceite.
			\item Falha de governança.
		\end{enumerate}
	\end{enumerate}
\end{frame}

\begin{frame}{Princípios do Gerenciamento de Projetos}
	\begin{enumerate}
		\setcounter{enumi}{66}
		\justifying
		\item Sobre Riscos Residuais:
		\begin{enumerate}[I]
			\item Permanecem após a resposta.
			\item Devem ser monitorados.
			\item Devem ser zerados sempre.
			\item Podem exigir planos de contingência.
		\end{enumerate}
		\begin{enumerate}[a]
			\item Se apenas I, II e IV estiverem corretas.
			\item Se apenas II e III estiverem corretas.
			\item \alert<2>{Se apenas I, II e IV estiverem corretas.}
			\item Se apenas IV estiver correta.
			\item Se apenas I e III estiverem corretas.
		\end{enumerate}
	\end{enumerate}
\end{frame}

\begin{frame}{Princípios do Gerenciamento de Projetos}
	\begin{enumerate}
		\setcounter{enumi}{67}
		\justifying
		\item A \rule{50pt}{0.75pt} permite ajustar-se a novas situações; a \rule{50pt}{0.75pt} garante voltar ao trilho após falhas.
		\begin{enumerate}[a]
			\item Resiliência / Adaptabilidade
			\item \alert<2>{Adaptabilidade / Resiliência}
			\item Flexibilidade / Rigidez
			\item Mudança / Estabilidade
			\item Velocidade / Precisão
		\end{enumerate}
	\end{enumerate}
\end{frame}

\begin{frame}{Princípios do Gerenciamento de Projetos}
	\begin{enumerate}
		\setcounter{enumi}{68}
		\justifying
		\item Favorecem adaptabilidade e resiliência:
		\begin{enumerate}[I]
			\item Transparência.
			\item Cultura de experimentação.
			\item Centralização total.
			\item Soluções modulares.
		\end{enumerate}
		\begin{enumerate}[a]
			\item Se apenas I, II e IV estiverem corretas.
			\item Se apenas II e IV estiverem corretas.
			\item \alert<2>{Se apenas I, II e IV estiverem corretas.}
			\item Se apenas IV estiver correta.
			\item Se apenas I e III estiverem corretas.
		\end{enumerate}
	\end{enumerate}
\end{frame}

\begin{frame}{Princípios do Gerenciamento de Projetos}
	\begin{enumerate}
		\setcounter{enumi}{69}
		\justifying
		\item Relacione conceitos à prática:
		\begin{columns}
			\begin{column}{0.3\textwidth}
				\begin{enumerate}[A]
					\item Adaptabilidade
					\item Resiliência
					\item Feedback
					\item Inspeção
				\end{enumerate}
			\end{column}
			\begin{column}{0.7\textwidth}
				\begin{itemize}
					\item Ajustar design após teste.
					\item Recuperar de atraso por desastre natural.
					\item Informação para validar direção.
					\item Verificação técnica de padrões.
				\end{itemize}
			\end{column}
		\end{columns}
		\begin{enumerate}[a]
			\item A, B, C, D.
			\item \alert<2>{B, A, C, D}
			\item C, D, A, B.
			\item D, C, B, A.
			\item A, C, B, D.
		\end{enumerate}
	\end{enumerate}
\end{frame}

\begin{frame}{Princípios do Gerenciamento de Projetos}
	\begin{enumerate}
		\setcounter{enumi}{70}
		\justifying
		\item Desenvolver resiliência perante problema:
		\begin{enumerate}[I]
			\item Aceitar a ocorrência.
			\item Estabilizar e avaliar danos.
			\item Encontrar caminhos alternativos.
			\item Extrair lições aprendidas.
		\end{enumerate}
		\begin{enumerate}[a]
			\item 4, 2, 1, 3.
			\item \alert<2>{1, 2, 3, 4}
			\item 2, 1, 4, 3.
			\item 3, 4, 1, 2.
			\item 1, 3, 2, 4.
		\end{enumerate}
	\end{enumerate}
\end{frame}

\begin{frame}{Princípios do Gerenciamento de Projetos}
	\begin{enumerate}
		\setcounter{enumi}{71}
		\justifying
		\item Abordagem adaptativa abraça mudança porque:
		\begin{enumerate}[a]
			\item Não exige planeamento.
			\item Custo é irrelevante.
			\item \alert<2>{Requisitos totais podem ser desconhecidos.}
			\item Projeto nunca termina.
			\item Não exige documentação.
		\end{enumerate}
	\end{enumerate}
\end{frame}

\begin{frame}{Princípios do Gerenciamento de Projetos}
	\begin{enumerate}
		\setcounter{enumi}{72}
		\justifying
		\item Equipe pivota projeto após lançamento de concorrente mais barato. Isso é:
		\begin{enumerate}[a]
			\item Falta de planeamento.
			\item \alert<2>{Alta adaptabilidade.}
			\item Desperdício.
			\item Incompetência.
			\item Falha de governança.
		\end{enumerate}
	\end{enumerate}
\end{frame}

\begin{frame}{Princípios do Gerenciamento de Projetos}
	\begin{enumerate}
		\setcounter{enumi}{73}
		\justifying
		\item Dificultam a resiliência:
		\begin{enumerate}[I]
			\item Punição a erros honestos.
			\item Burocracia lenta.
			\item Silos de informação.
			\item Autonomia delegada.
		\end{enumerate}
		\begin{enumerate}[a]
			\item Se apenas I, II e III estiverem corretas.
			\item Se apenas II e IV estiverem corretas.
			\item \alert<2>{Se apenas I, II e III estiverem corretas.}
			\item Se apenas IV estiver correta.
			\item Se apenas I e III estiverem corretas.
		\end{enumerate}
	\end{enumerate}
\end{frame}

\begin{frame}{Princípios do Gerenciamento de Projetos}
	\begin{enumerate}
		\setcounter{enumi}{74}
		\justifying
		\item Papel do projeto na organização:
		\begin{enumerate}[a]
			\item Manter status quo.
			\item \alert<2>{Veículo para levar ao estado futuro.}
			\item Substituir pessoas por máquinas.
			\item Criar burocracia.
			\item Impor tecnologia.
		\end{enumerate}
	\end{enumerate}
\end{frame}

\begin{frame}{Princípios do Gerenciamento de Projetos}
	\begin{enumerate}
		\setcounter{enumi}{75}
		\justifying
		\item Focar na \rule{50px}{0.7pt} dos novos resultados e no \rule{50px}{0.7pt} dos benefícios.
		\begin{enumerate}[a]
			\item Adoção / Alcance (ou Realização)
			\item \alert<2>{Adoção / Alcance (ou Realização)}
			\item Entrega / Medição
			\item Execução / Finalização
			\item Aceitação / Transferência
		\end{enumerate}
	\end{enumerate}
\end{frame}

\begin{frame}{Princípios do Gerenciamento de Projetos}
	\begin{enumerate}
		\setcounter{enumi}{76}
		\justifying
		\item Fundamentais para a mudança:
		\begin{enumerate}[I]
			\item Comunicação da visão.
			\item Treino e suporte.
			\item Mitigar resistência.
			\item Força e punição.
		\end{enumerate}
		\begin{enumerate}[a]
			\item Se apenas I, II e III estiverem corretas.
			\item Se apenas II e IV estiverem corretas.
			\item \alert<2>{Se apenas I, II e III estiverem corretas.}
			\item Se apenas IV estiver correta.
			\item Se todas estiverem corretas.
		\end{enumerate}
	\end{enumerate}
\end{frame}

\begin{frame}{Princípios do Gerenciamento de Projetos}
	\begin{enumerate}
		\setcounter{enumi}{77}
		\justifying
		\item Relacione conceitos de mudança:
		\begin{columns}
			\begin{column}{0.3\textwidth}
				\begin{enumerate}[A]
					\item Estado Atual
					\item Estado Futuro
					\item Transição
					\item Sustentação
				\end{enumerate}
			\end{column}
			\begin{column}{0.7\textwidth}
				\begin{itemize}
					\item Onde a organização está hoje.
					\item O objetivo pretendido.
					\item Período de aprendizagem e trabalho.
					\item Garantir que a mudança 'grude'.
				\end{itemize}
			\end{column}
		\end{columns}
		\begin{enumerate}[78]
			\item A, B, C, D.
			\item \alert<2>{A, C, D, B}
			\item C, D, A, B.
			\item D, C, B, A.
			\item B, A, D, C.
		\end{enumerate}
	\end{enumerate}
\end{frame}

\begin{frame}{Princípios do Gerenciamento de Projetos}
	\begin{enumerate}
		\setcounter{enumi}{79}
		\justifying
		\item Facilitar transição de sistema:
		\begin{enumerate}[I]
			\item Comunicar necessidade.
			\item Treinar e realizar pilotos.
			\item Entregar e desativar antigo.
			\item Reforçar adoção.
		\end{enumerate}
		\begin{enumerate}[a]
			\item 4, 2, 1, 3.
			\item \alert<2>{1, 2, 3, 4}
			\item 2, 1, 4, 3.
			\item 3, 4, 1, 2.
			\item 1, 3, 2, 4.
		\end{enumerate}
	\end{enumerate}
\end{frame}

\begin{frame}{Princípios do Gerenciamento de Projetos}
	\begin{enumerate}
		\setcounter{enumi}{80}
		\justifying
		\item Mudança bem-sucedida ocorre quando:
		\begin{enumerate}[a]
			\item Orçamento é seguido à risca.
			\item \alert<2>{Usuários adotam solução e benefícios aparecem.}
			\item Gerente sai logo.
			\item Ninguém nota mudança.
			\item Projeto é cancelado cedo.
		\end{enumerate}
	\end{enumerate}
\end{frame}

\begin{frame}{Princípios do Gerenciamento de Projetos}
	\begin{enumerate}
		\setcounter{enumi}{81}
		\justifying
		\item ERP funciona mas funcionários usam folha de cálculo externa. Falha em:
		\begin{enumerate}[a]
			\item Navegar na complexidade tecnológica.
			\item Focar na qualidade do código.
			\item \alert<2>{Possibilitar a mudança.}
			\item Cronograma de hardware.
			\item Normas de segurança.
		\end{enumerate}
	\end{enumerate}
\end{frame}

\begin{frame}{Princípios do Gerenciamento de Projetos}
	\begin{enumerate}
		\setcounter{enumi}{82}
		\justifying
		\item Responsáveis por possibilitar a mudança:
		\begin{enumerate}[I]
			\item Liderança.
			\item Equipa.
			\item Gerentes funcionais.
			\item Agentes de mudança.
		\end{enumerate}
		\begin{enumerate}[a]
			\item Se todas estiverem corretas.
			\item Se apenas II e IV estiverem corretas.
			\item \alert<2>{Se todas estiverem corretas.}
			\item Se apenas IV estiver correta.
			\item Se apenas I e III estiverem corretas.
		\end{enumerate}
	\end{enumerate}
\end{frame}
%===============================================================================
% FIM ==========================================================================
%===============================================================================

\begin{frame}
	\centering
	\Large
	FIM
\end{frame}

\end{document}